\section{Discussion}
\label{sec:discussion}

\paragraph{When to escalate.} Average success, action divergence and replacement gain answer different questions and can disagree. The four W4 cells have small end-to-end differences with wide intervals (Appendix~\ref{app:e2e}). Their Bank P states replace actions with 15.5--20.9\% turnover, but the pooled replacement gains are near zero. % src: analysis/fork/out/dev_all_20261001T070326Z.md
Reporting only the first two cannot tell whether a BF16 fallback is worth its cost. Text actions should be evaluated as actions, but, since their outcome labels did not reproduce reliably under another simulator (\S\ref{sec:r3}), not reported as a property of the action without a simulator control. We recommend reporting $r$, $d$ and $M$ next to any net gain, each with its denominator, and keeping fixed-pool and new-task intervals apart.

\paragraph{Which part to replace.} The best unit of replacement is not universal. On the fixed task pool, the q8 W4 cells favour replacing the current action and a severe W3 cell the continuation (\S\ref{sec:r2}), but these orderings are exploratory and unconfirmed under a second simulator (\S\ref{sec:r3}). Measure $G$, $H$ and $F$ per configuration and environment.

\paragraph{Escalation can hurt.} Replacements damage as well as rescue (\S\ref{sec:r1}); B continuing lowers success in q4 KV8 (\S\ref{sec:r2}), and on unseen telecom tasks action replacement lowers it by 4.6\pp{} (\S\ref{sec:r3}). % src: analysis/fork/out/holdout_confirm_20261002T012353Z.md

